\documentclass[11pt]{article}
\usepackage{preamble}
\setlength{\parskip}{4pt}

\begin{document}

\daybanner{AP Statistics \textperiodcentered\ Unit 3 \textperiodcentered\ Topic 3.3 \textperiodcentered\ TEACHER KEY}{Six Activations in Twelve Users}

\sectionbanner{CED TETHER}

{\small
\begin{itemize}[leftmargin=1.5em,itemsep=2pt,topsep=2pt]
  \item \textbf{VAR-1.H} Identify questions suggested by variation in the shapes of distributions of samples taken from the same population. \textbf{[Skill 1.A]}
  \item \textbf{VAR-1.H.1} Variation in shapes of data distributions may be random or not.
\end{itemize}
}
\textbf{Essential question:} When should a simulated sampling distribution be trusted more than a normal approximation?\par
\textbf{DOK-3 skill:} 4.B \textperiodcentered\ \textbf{FRQ pattern:} {\small\texttt{normal-\allowbreak{}approximation-\allowbreak{}vs-\allowbreak{}simulation-\allowbreak{}for-\allowbreak{}a-\allowbreak{}skewed-\allowbreak{}tail}}\par
\textbf{DOK rationale:} Part (c) coordinates an expected-count check, the simulated distribution's shape, and two tail estimates to choose a defensible method, bound the evidence, and explain the consequence of the rejected analysis.\par

\textbf{Self-paced bonus sheet.} Students take it when they choose and turn it in when they finish. Everything they need is printed on the sheet. Score part (c) only, E / P / I, by hand --- never AI-graded or auto-scored. (a) and (b) are the ladder up; use them to see where a thin (c) came from.\par

\textbf{Questions to ask}
\begin{itemize}[leftmargin=1.5em,itemsep=2pt,topsep=2pt]
  \item What does one bar in this simulated distribution count?
  \item Which expected count warns you that a symmetric curve may miss the tail?
  \item Why do simulated samples with more than six activations still count?
  \item Can the two methods disagree in size yet lead to the same 5 percent decision?
\end{itemize}

\begin{callout}{\IconWarn\ LOOK FOR}{calloutred}
\begin{itemize}[leftmargin=1.5em,itemsep=2pt,topsep=2pt]
  \item Treating every sampling distribution as normal regardless of expected counts.
  \item Counting exactly six activations instead of six or more.
  \item Saying a small chance probability proves why the activation rate differs.
\end{itemize}
\end{callout}

\sectionbanner{SCORING PART (c) --- E / P / I}

\textbf{Expected elements}
\begin{itemize}[leftmargin=1.5em,itemsep=2pt,topsep=2pt]
  \item \textbf{judgment-1} (required) --- Chooses the simulation using failed expected counts 2.4 and 9.6 or the right-skewed shape.
  \item \textbf{judgment-2} (required) --- Uses 4/200=0.020 below 0.05 as evidence against the model, without claiming proof.
  \item \textbf{judgment-3} (required) --- Explains that 0.005 makes the result seem four times rarer and overstates the evidence.
\end{itemize}

{\small\renewcommand{\arraystretch}{1.25}
\noindent\begin{tabular}{|p{0.5in}|p{5.9in}|}
\hline
\textbf{E} & Meets all three checks with correct contextual reasoning; equivalent wording is acceptable. \\ \hline
\textbf{P} & Shows at least one substantive correct check, but has an omission or error that prevents E. \\ \hline
\textbf{I} & Shows none of the three checks with substantive correct reasoning; a choice alone is insufficient. \\ \hline
\end{tabular}}

\textbf{Common mistakes}
\begin{itemize}[leftmargin=1.5em,itemsep=2pt,topsep=2pt]
  \item Using $6/12=0.50$ as the chance probability
  \item Counting only the bar at 6 without recognizing that all larger results would also be at least as extreme
  \item Checking observed successes and failures instead of expected counts under $p_0$ for this model
  \item Calling the two percent estimate proof that the developer intentionally misstated the rate
\end{itemize}

\clearpage
\focusbanner{THE PROBLEM --- annotated}

Suppose a developer's model says $p_0=0.20$ of new app users activate voice commands. An SRS of 12 users includes 6 activations, so $\hat p=0.50$. The chart shows 200 simulated samples of 12 independent users under this model.\par\smallskip \textbf{Mara:} ``A normal model for $\hat p$ gives $P(\hat p\ge 0.50)\approx0.005$, so that is the chance probability I would report.''\par \textbf{Jules:} ``The simulated distribution is not very normal. I would estimate the chance probability directly from its tail.''\par
\begin{center}
\begin{tikzpicture}
\begin{axis}[
  width=0.98\linewidth,height=1.75in,
  ybar,bar width=10pt,xmin=-0.6,xmax=6.6,ymin=0,ymax=70,
  xtick={0,1,2,3,4,5,6},ytick={0,20,40,60},
  xlabel={Activations in a sample of 12},ylabel={Simulated samples},
  tick label style={font=\small},label style={font=\small},
  nodes near coords,nodes near coords style={font=\scriptsize},
  axis lines*=left
]
\addplot[draw=dayblue,fill=dayblue!20] coordinates {(0,14) (1,40) (2,58) (3,47) (4,25) (5,12) (6,4)};
\end{axis}
\end{tikzpicture}
\end{center}
\begin{firsttakebox}{FIRST TAKE --- one sentence before you work the parts}
Whose method would you trust more, Mara's smooth curve or Jules's repeated trials? Give one reason from the picture.\par
\answer{Not graded. Preserve the normal-model versus matched-simulation tension.}
\end{firsttakebox}

\begin{callout}{\IconBook\ CHECK THE SHAPE BEFORE USING NORMAL (keep on the page)}{calloutgreen}
Under a proposed proportion $p_0$, check expected counts $np_0\geq10$ and $n(1-p_0)\geq10$ before using a normal approximation.
A matched simulation estimates the chance probability by the number of simulated results at least as large as the observed result, divided by the total simulations.
For this task, a chance probability below $5\%$ counts as convincing evidence against the proposed model.
Evidence against a model does not prove a particular cause.
\end{callout}

\dokbadge{1}~\textbf{(a)} State the most common number of activations in the simulated samples and describe the distribution as roughly symmetric, skewed left, or skewed right.\par
\answer{Two activations is most common (58 simulations). The distribution is skewed right.}

\dokbadge{2}~\textbf{(b)} Use the chart to estimate $P(\hat p\ge0.50)$ under $p_0=0.20$. Then calculate $np_0$ and $n(1-p_0)$ and assess the expected-count condition for a normal approximation.\par
\answer{The simulated tail is $4/200=0.020$. Expected counts $12(0.20)=2.4$ and $12(0.80)=9.6$ both fall below 10, so the normal condition fails.}

\dokbadge{3}~\textbf{(c)} In 3--4 sentences, choose a method using the simulated tail and either the shape or expected counts. Apply the $5\%$ evidence rule. Explain how reporting 0.005 instead of 0.020 would change a reader's impression.\par
\answer{Use Jules's simulation: the failed expected counts and right skew weaken the normal approximation. The tail $0.020<0.05$ gives evidence against the model, without proving a cause. Reporting $0.005$ makes the result seem four times rarer and overstates the evidence.}

\end{document}
